\documentclass{dlproblemset}
\usepackage{tikz}
\usetikzlibrary{positioning}
\tikzset{every picture/.style={line width=0.75pt}}

% ---- Histogramas de pré-ativações usados na questão das quatro inicializações ----
% Pré-ativações de uma MLP ReLU simulada (n = 256, minilote de 1000 instâncias).
% Eixo horizontal comum em [-8, 8]; altura normalizada pelo pico de cada painel.
\newcommand{\histpanel}[1]{%
\begin{tikzpicture}[x=0.19cm, y=1.1cm, baseline=(current bounding box.center)]
  \foreach \gx in {-4,0,4} \draw[gray!25] (\gx,0) -- (\gx,1.05);
  \fill[black!22] plot[const plot] coordinates {#1} -- (8,0) -- (-8,0) -- cycle;
  \draw[black!75, thick] plot[const plot] coordinates {#1};
  \draw[gray!70] (-8,0) rectangle (8,1.05);
  \foreach \t in {-8,-4,0,4,8} \node[font=\tiny, below, inner sep=1.5pt] at (\t,0) {$\t$};
\end{tikzpicture}}

\title{Inicialização de Pesos}
\subtitle{Lista de Exercícios}

\begin{document}

\begin{enumerate}[1)]

    \item Você inicializa todos os pesos de um MLP com um mesmo valor constante $c \neq 0$ e todos os bias com zero. As camadas ocultas usam a função de ativação tanh. Qual o principal problema que esse esquema de inicialização causa durante o treinamento?

    \begin{enumerate}[(a)]
        \item Os gradientes da primeira camada são nulos já na primeira iteração e os pesos dessa camada nunca são atualizados.
        \item A função de custo se torna não diferenciável nos pontos em que pesos de camadas distintas possuem o mesmo valor.
        \item As unidades de uma mesma camada computam a mesma saída e recebem gradientes idênticos, permanecendo cópias umas das outras.
        \item As pré-ativações crescem exponencialmente com a profundidade já no primeiro \textit{forward pass}, causando \textit{overflow} numérico.
    \end{enumerate}
    %c

    \item Um MLP profundo com ativações sigmoid nas camadas ocultas tem os pesos inicializados com $\theta \sim \mathcal{N}(0, \sigma^2)$, onde $\sigma$ é muito maior do que o recomendado pelas heurísticas de inicialização (He, Xavier). Qual comportamento você espera observar nas primeiras iterações do treinamento?

    \begin{enumerate}[(a)]
        \item Unidades mortas em todas as camadas ocultas, com saída exatamente zero independentemente da entrada apresentada.
        \item Pré-ativações com magnitude grande, unidades saturadas e gradientes próximos de zero nas camadas iniciais da rede.
        \item Pré-ativações próximas de zero, unidades operando na região aproximadamente linear e gradientes estáveis em todas as camadas.
        \item Convergência acelerada nas primeiras épocas, seguida de estagnação quando a magnitude dos pesos diminui ao longo do treino.
    \end{enumerate}
    %b

    \item Que problemas podemos ter com unidades na camada oculta que usam a função sigmoid?

    \item Que problemas podemos ter com unidades na camada oculta que usam a função relu?

    \item Qual a motivação das inicializações \textit{He} e \textit{Xavier}.

    \item Seja $a$ uma variável aleatória contínua, com distribuição simétrica em torno de $E[a]=0$ e variância $\mathrm{Var}[a]=\sigma^2$. Mostre que
    \begin{equation*}
        E\!\left[\mathrm{ReLU}(a)^2\right]=\frac{\sigma^2}{2}
    \end{equation*}
    \textit{Dica: escreva a esperança como uma integral, use a simetria da distribuição para relacionar a integral na metade positiva do domínio com $E[a^2]$ e lembre que $E[a^2]=\mathrm{Var}[a]+(E[a])^2$.}

    \item Considere a pré-ativação da unidade $i$ na camada $l$ de um MLP cuja camada anterior possui $n^{(l-1)}$ unidades:
    \begin{equation*}
        z^{(l)}_i = b^{(l)}_i + \sum_{j=1}^{n^{(l-1)}}\theta^{(l)}_{ij}\,\varphi\!\left(z^{(l-1)}_j\right)
    \end{equation*}
    \noindent onde os bias são inicializados em zero, os pesos $\theta^{(l)}_{ij}\sim\mathcal{N}(0,\sigma_\theta^2)$ são independentes entre si e independentes de $\varphi(z^{(l-1)}_j)$, e as pré-ativações $z^{(l-1)}_j$ têm distribuição simétrica em torno de zero com variância $\sigma^2_{z^{(l-1)}}$. Resolva os itens abaixo:
    \begin{enumerate}[a)]
        \item Mostre que $E[z^{(l)}_i]=0$.
        \item Supondo $\varphi=\mathrm{ReLU}$, use o resultado do exercício anterior para mostrar a recorrência $\sigma^2_{z^{(l)}} = \frac{n^{(l-1)}}{2}\,\sigma_\theta^2\,\sigma^2_{z^{(l-1)}}$.
        \item Imponha que a variância se preserve de uma camada para a outra, isto é, $\sigma^2_{z^{(l)}}=\sigma^2_{z^{(l-1)}}$, e derive a variância da inicialização \textit{He}.
        \item Supondo agora $\varphi=\tanh$ e usando a aproximação $\varphi(z)\approx z$ para $z$ próximo da origem, refaça o argumento dos itens (b) e (c) e derive a variância da inicialização \textit{Xavier}.
    \end{enumerate}

    \item Sobre a inicialização dos pesos de uma rede neural profunda com função de ativação ReLU, qual afirmação é correta?
    \begin{enumerate}[(a)]
        \item Inicializar todos os pesos com zero é adequado para ReLU; o gradiente assimétrico da ReLU faz com que diferentes unidades aprendam características distintas mesmo partindo de pesos idênticos.
        \item Inicializar todos os pesos com a mesma constante não-nula é adequado, desde que a constante seja calibrada em função do número de camadas.
        \item A inicialização de Xavier (variância $1/n_{\text{in}}$ ou $2/(n_{\text{in}} + n_{\text{out}})$) é a recomendada para ReLU, pois assume que a função de ativação se comporta de forma aproximadamente linear.
        \item A inicialização de He (variância $2/n_{\text{in}}$) é recomendada para ReLU, compensando o fato de cerca de metade das pré-ativações serem zeradas.
        \item Qualquer inicialização aleatória uniforme em $[-1, 1]$ funciona igualmente bem, independentemente da profundidade da rede e da função de ativação utilizada.
        \item Para ReLU, recomenda-se uma inicialização determinística baseada na projeção PCA do conjunto de treinamento sobre os pesos da primeira camada.
        \item A inicialização aleatória dos pesos é desnecessária; basta inicializar os vieses (\textit{biases}) com valores aleatórios para quebrar a simetria.
        \item Pesos amostrados de $\mathcal{N}(0, 1)$ (variância 1, sem reescala) são apropriados para qualquer arquitetura com ReLU, desde que a taxa de aprendizado seja suficientemente pequena.
    \end{enumerate}
    %d

    \item Um MLP tem entrada de dimensão $256$ e $4$ camadas ocultas com $n = 256$ unidades ReLU cada. As componentes da entrada são independentes, com distribuição $\mathcal{N}(0, 1)$, os vieses são inicializados em zero e os pesos com $\theta^{(l)}_{ij} \sim \mathcal{N}(0, \sigma_\theta^2)$. A rede foi inicializada com quatro valores diferentes de $\sigma_\theta^2$ e, para cada um, mediu-se o histograma das pré-ativações $\mathbf{Z}^{(1)}, \ldots, \mathbf{Z}^{(4)}$ em um minilote de $1000$ instâncias. Cada linha (P, Q, R e S) corresponde a uma inicialização. O eixo horizontal é o mesmo em todos os painéis, e a altura de cada painel está normalizada pelo seu próprio pico.

    \begin{center}
    \setlength{\tabcolsep}{3pt}
    \begin{tabular}{c@{\hspace{6pt}}cccc}
     & $\mathbf{Z}^{(1)}$ & $\mathbf{Z}^{(2)}$ & $\mathbf{Z}^{(3)}$ & $\mathbf{Z}^{(4)}$ \\
    \textbf{P} & \histpanel{(-8.0,0.0000) (-7.8,0.0000) (-7.6,0.0000) (-7.4,0.0000) (-7.2,0.0000) (-7.0,0.0000) (-6.8,0.0000) (-6.6,0.0000) (-6.4,0.0000) (-6.2,0.0000) (-6.0,0.0000) (-5.8,0.0000) (-5.6,0.0000) (-5.4,0.0000) (-5.2,0.0000) (-5.0,0.0000) (-4.8,0.0000) (-4.6,0.0000) (-4.4,0.0002) (-4.2,0.0003) (-4.0,0.0004) (-3.8,0.0010) (-3.6,0.0025) (-3.4,0.0049) (-3.2,0.0080) (-3.0,0.0154) (-2.8,0.0286) (-2.6,0.0483) (-2.4,0.0750) (-2.2,0.1118) (-2.0,0.1662) (-1.8,0.2334) (-1.6,0.3302) (-1.4,0.4315) (-1.2,0.5589) (-1.0,0.6683) (-0.8,0.7911) (-0.6,0.8984) (-0.4,0.9539) (-0.2,1.0000) (0.0,0.9889) (0.2,0.9646) (0.4,0.8935) (0.6,0.7910) (0.8,0.6685) (1.0,0.5409) (1.2,0.4260) (1.4,0.3307) (1.6,0.2295) (1.8,0.1693) (2.0,0.1123) (2.2,0.0742) (2.4,0.0458) (2.6,0.0265) (2.8,0.0168) (3.0,0.0083) (3.2,0.0053) (3.4,0.0027) (3.6,0.0014) (3.8,0.0008) (4.0,0.0002) (4.2,0.0000) (4.4,0.0001) (4.6,0.0000) (4.8,0.0000) (5.0,0.0000) (5.2,0.0000) (5.4,0.0000) (5.6,0.0000) (5.8,0.0000) (6.0,0.0000) (6.2,0.0000) (6.4,0.0000) (6.6,0.0000) (6.8,0.0000) (7.0,0.0000) (7.2,0.0000) (7.4,0.0000) (7.6,0.0000) (7.8,0.0000) (8.0,0)} & \histpanel{(-8.0,0.0000) (-7.8,0.0000) (-7.6,0.0000) (-7.4,0.0000) (-7.2,0.0000) (-7.0,0.0000) (-6.8,0.0000) (-6.6,0.0000) (-6.4,0.0000) (-6.2,0.0000) (-6.0,0.0000) (-5.8,0.0000) (-5.6,0.0000) (-5.4,0.0000) (-5.2,0.0000) (-5.0,0.0000) (-4.8,0.0000) (-4.6,0.0000) (-4.4,0.0000) (-4.2,0.0000) (-4.0,0.0000) (-3.8,0.0000) (-3.6,0.0000) (-3.4,0.0001) (-3.2,0.0001) (-3.0,0.0004) (-2.8,0.0011) (-2.6,0.0029) (-2.4,0.0066) (-2.2,0.0141) (-2.0,0.0313) (-1.8,0.0622) (-1.6,0.1152) (-1.4,0.1978) (-1.2,0.3102) (-1.0,0.4550) (-0.8,0.6305) (-0.6,0.8029) (-0.4,0.9377) (-0.2,1.0000) (0.0,0.9893) (0.2,0.8887) (0.4,0.7502) (0.6,0.5812) (0.8,0.4017) (1.0,0.2626) (1.2,0.1576) (1.4,0.0863) (1.6,0.0452) (1.8,0.0214) (2.0,0.0095) (2.2,0.0040) (2.4,0.0019) (2.6,0.0005) (2.8,0.0004) (3.0,0.0002) (3.2,0.0000) (3.4,0.0000) (3.6,0.0000) (3.8,0.0000) (4.0,0.0000) (4.2,0.0000) (4.4,0.0000) (4.6,0.0000) (4.8,0.0000) (5.0,0.0000) (5.2,0.0000) (5.4,0.0000) (5.6,0.0000) (5.8,0.0000) (6.0,0.0000) (6.2,0.0000) (6.4,0.0000) (6.6,0.0000) (6.8,0.0000) (7.0,0.0000) (7.2,0.0000) (7.4,0.0000) (7.6,0.0000) (7.8,0.0000) (8.0,0)} & \histpanel{(-8.0,0.0000) (-7.8,0.0000) (-7.6,0.0000) (-7.4,0.0000) (-7.2,0.0000) (-7.0,0.0000) (-6.8,0.0000) (-6.6,0.0000) (-6.4,0.0000) (-6.2,0.0000) (-6.0,0.0000) (-5.8,0.0000) (-5.6,0.0000) (-5.4,0.0000) (-5.2,0.0000) (-5.0,0.0000) (-4.8,0.0000) (-4.6,0.0000) (-4.4,0.0000) (-4.2,0.0000) (-4.0,0.0000) (-3.8,0.0000) (-3.6,0.0000) (-3.4,0.0000) (-3.2,0.0000) (-3.0,0.0000) (-2.8,0.0000) (-2.6,0.0000) (-2.4,0.0000) (-2.2,0.0002) (-2.0,0.0006) (-1.8,0.0023) (-1.6,0.0072) (-1.4,0.0202) (-1.2,0.0549) (-1.0,0.1355) (-0.8,0.2918) (-0.6,0.5443) (-0.4,0.8085) (-0.2,1.0000) (0.0,0.9973) (0.2,0.8191) (0.4,0.5652) (0.6,0.3119) (0.8,0.1454) (1.0,0.0596) (1.2,0.0214) (1.4,0.0075) (1.6,0.0018) (1.8,0.0004) (2.0,0.0001) (2.2,0.0000) (2.4,0.0000) (2.6,0.0000) (2.8,0.0000) (3.0,0.0000) (3.2,0.0000) (3.4,0.0000) (3.6,0.0000) (3.8,0.0000) (4.0,0.0000) (4.2,0.0000) (4.4,0.0000) (4.6,0.0000) (4.8,0.0000) (5.0,0.0000) (5.2,0.0000) (5.4,0.0000) (5.6,0.0000) (5.8,0.0000) (6.0,0.0000) (6.2,0.0000) (6.4,0.0000) (6.6,0.0000) (6.8,0.0000) (7.0,0.0000) (7.2,0.0000) (7.4,0.0000) (7.6,0.0000) (7.8,0.0000) (8.0,0)} & \histpanel{(-8.0,0.0000) (-7.8,0.0000) (-7.6,0.0000) (-7.4,0.0000) (-7.2,0.0000) (-7.0,0.0000) (-6.8,0.0000) (-6.6,0.0000) (-6.4,0.0000) (-6.2,0.0000) (-6.0,0.0000) (-5.8,0.0000) (-5.6,0.0000) (-5.4,0.0000) (-5.2,0.0000) (-5.0,0.0000) (-4.8,0.0000) (-4.6,0.0000) (-4.4,0.0000) (-4.2,0.0000) (-4.0,0.0000) (-3.8,0.0000) (-3.6,0.0000) (-3.4,0.0000) (-3.2,0.0000) (-3.0,0.0000) (-2.8,0.0000) (-2.6,0.0000) (-2.4,0.0000) (-2.2,0.0000) (-2.0,0.0000) (-1.8,0.0000) (-1.6,0.0001) (-1.4,0.0013) (-1.2,0.0085) (-1.0,0.0403) (-0.8,0.1341) (-0.6,0.3412) (-0.4,0.6837) (-0.2,1.0000) (0.0,0.9993) (0.2,0.6824) (0.4,0.3201) (0.6,0.1117) (0.8,0.0283) (1.0,0.0054) (1.2,0.0006) (1.4,0.0001) (1.6,0.0000) (1.8,0.0000) (2.0,0.0000) (2.2,0.0000) (2.4,0.0000) (2.6,0.0000) (2.8,0.0000) (3.0,0.0000) (3.2,0.0000) (3.4,0.0000) (3.6,0.0000) (3.8,0.0000) (4.0,0.0000) (4.2,0.0000) (4.4,0.0000) (4.6,0.0000) (4.8,0.0000) (5.0,0.0000) (5.2,0.0000) (5.4,0.0000) (5.6,0.0000) (5.8,0.0000) (6.0,0.0000) (6.2,0.0000) (6.4,0.0000) (6.6,0.0000) (6.8,0.0000) (7.0,0.0000) (7.2,0.0000) (7.4,0.0000) (7.6,0.0000) (7.8,0.0000) (8.0,0)} \\[0.4em]
    \textbf{Q} & \histpanel{(-8.0,0.0004) (-7.8,0.0008) (-7.6,0.0003) (-7.4,0.0010) (-7.2,0.0021) (-7.0,0.0027) (-6.8,0.0041) (-6.6,0.0040) (-6.4,0.0073) (-6.2,0.0103) (-6.0,0.0125) (-5.8,0.0162) (-5.6,0.0228) (-5.4,0.0282) (-5.2,0.0388) (-5.0,0.0479) (-4.8,0.0597) (-4.6,0.0763) (-4.4,0.0961) (-4.2,0.1188) (-4.0,0.1496) (-3.8,0.1776) (-3.6,0.2132) (-3.4,0.2452) (-3.2,0.2967) (-3.0,0.3533) (-2.8,0.4026) (-2.6,0.4559) (-2.4,0.5169) (-2.2,0.5604) (-2.0,0.6292) (-1.8,0.6916) (-1.6,0.7478) (-1.4,0.8241) (-1.2,0.8594) (-1.0,0.9059) (-0.8,0.9593) (-0.6,0.9930) (-0.4,0.9950) (-0.2,0.9888) (0.0,1.0000) (0.2,0.9924) (0.4,0.9851) (0.6,0.9430) (0.8,0.9168) (1.0,0.8490) (1.2,0.8132) (1.4,0.7601) (1.6,0.6936) (1.8,0.6366) (2.0,0.5755) (2.2,0.5081) (2.4,0.4543) (2.6,0.3941) (2.8,0.3440) (3.0,0.2995) (3.2,0.2577) (3.4,0.2085) (3.6,0.1794) (3.8,0.1485) (4.0,0.1214) (4.2,0.1030) (4.4,0.0707) (4.6,0.0635) (4.8,0.0471) (5.0,0.0420) (5.2,0.0282) (5.4,0.0250) (5.6,0.0179) (5.8,0.0131) (6.0,0.0106) (6.2,0.0073) (6.4,0.0039) (6.6,0.0031) (6.8,0.0027) (7.0,0.0018) (7.2,0.0016) (7.4,0.0007) (7.6,0.0010) (7.8,0.0001) (8.0,0)} & \histpanel{(-8.0,0.0207) (-7.8,0.0277) (-7.6,0.0260) (-7.4,0.0380) (-7.2,0.0399) (-7.0,0.0467) (-6.8,0.0588) (-6.6,0.0644) (-6.4,0.0817) (-6.2,0.0909) (-6.0,0.0999) (-5.8,0.1206) (-5.6,0.1377) (-5.4,0.1550) (-5.2,0.1718) (-5.0,0.1984) (-4.8,0.2264) (-4.6,0.2550) (-4.4,0.2801) (-4.2,0.3219) (-4.0,0.3681) (-3.8,0.3852) (-3.6,0.4293) (-3.4,0.4774) (-3.2,0.5005) (-3.0,0.5459) (-2.8,0.5835) (-2.6,0.6350) (-2.4,0.6702) (-2.2,0.7243) (-2.0,0.7582) (-1.8,0.8047) (-1.6,0.8490) (-1.4,0.8652) (-1.2,0.8985) (-1.0,0.9313) (-0.8,0.9394) (-0.6,0.9567) (-0.4,0.9827) (-0.2,0.9849) (0.0,1.0000) (0.2,0.9705) (0.4,0.9684) (0.6,0.9433) (0.8,0.9262) (1.0,0.9102) (1.2,0.8848) (1.4,0.8491) (1.6,0.8167) (1.8,0.7732) (2.0,0.7167) (2.2,0.7045) (2.4,0.6554) (2.6,0.6096) (2.8,0.5750) (3.0,0.5166) (3.2,0.4954) (3.4,0.4493) (3.6,0.4027) (3.8,0.3684) (4.0,0.3333) (4.2,0.3040) (4.4,0.2593) (4.6,0.2399) (4.8,0.2049) (5.0,0.1756) (5.2,0.1607) (5.4,0.1454) (5.6,0.1121) (5.8,0.1101) (6.0,0.0876) (6.2,0.0730) (6.4,0.0675) (6.6,0.0523) (6.8,0.0467) (7.0,0.0393) (7.2,0.0316) (7.4,0.0309) (7.6,0.0216) (7.8,0.0190) (8.0,0)} & \histpanel{(-8.0,0.1451) (-7.8,0.1544) (-7.6,0.1620) (-7.4,0.1762) (-7.2,0.2013) (-7.0,0.2163) (-6.8,0.2403) (-6.6,0.2580) (-6.4,0.2833) (-6.2,0.2918) (-6.0,0.3051) (-5.8,0.3530) (-5.6,0.3724) (-5.4,0.4030) (-5.2,0.4234) (-5.0,0.4608) (-4.8,0.4964) (-4.6,0.5131) (-4.4,0.5342) (-4.2,0.5549) (-4.0,0.6025) (-3.8,0.6337) (-3.6,0.6601) (-3.4,0.6894) (-3.2,0.7196) (-3.0,0.7433) (-2.8,0.7587) (-2.6,0.8266) (-2.4,0.8209) (-2.2,0.8515) (-2.0,0.8624) (-1.8,0.9040) (-1.6,0.8985) (-1.4,0.9348) (-1.2,0.9570) (-1.0,0.9618) (-0.8,0.9800) (-0.6,0.9641) (-0.4,1.0000) (-0.2,0.9842) (0.0,0.9740) (0.2,0.9643) (0.4,0.9532) (0.6,0.9810) (0.8,0.9601) (1.0,0.9382) (1.2,0.9363) (1.4,0.9038) (1.6,0.8899) (1.8,0.8568) (2.0,0.8260) (2.2,0.8101) (2.4,0.8011) (2.6,0.7656) (2.8,0.7477) (3.0,0.6966) (3.2,0.6865) (3.4,0.6667) (3.6,0.6468) (3.8,0.5924) (4.0,0.5709) (4.2,0.5422) (4.4,0.5171) (4.6,0.4776) (4.8,0.4479) (5.0,0.4175) (5.2,0.3726) (5.4,0.3719) (5.6,0.3348) (5.8,0.3072) (6.0,0.2913) (6.2,0.2766) (6.4,0.2439) (6.6,0.2340) (6.8,0.2169) (7.0,0.2063) (7.2,0.1768) (7.4,0.1675) (7.6,0.1544) (7.8,0.1308) (8.0,0)} & \histpanel{(-8.0,0.2992) (-7.8,0.3500) (-7.6,0.3401) (-7.4,0.3920) (-7.2,0.3887) (-7.0,0.4218) (-6.8,0.4205) (-6.6,0.4345) (-6.4,0.4744) (-6.2,0.5075) (-6.0,0.5088) (-5.8,0.5182) (-5.6,0.5630) (-5.4,0.5825) (-5.2,0.6026) (-5.0,0.6258) (-4.8,0.6372) (-4.6,0.6660) (-4.4,0.7192) (-4.2,0.7024) (-4.0,0.7278) (-3.8,0.7235) (-3.6,0.7777) (-3.4,0.7937) (-3.2,0.7856) (-3.0,0.8067) (-2.8,0.8166) (-2.6,0.8731) (-2.4,0.8649) (-2.2,0.9125) (-2.0,0.9361) (-1.8,0.9191) (-1.6,0.9570) (-1.4,0.9423) (-1.2,0.9316) (-1.0,0.9415) (-0.8,0.9679) (-0.6,0.9819) (-0.4,0.9791) (-0.2,1.0000) (0.0,0.9763) (0.2,0.9845) (0.4,0.9921) (0.6,0.9789) (0.8,0.9524) (1.0,0.9506) (1.2,0.9359) (1.4,0.9428) (1.6,0.9303) (1.8,0.9387) (2.0,0.8949) (2.2,0.8692) (2.4,0.8603) (2.6,0.8408) (2.8,0.8293) (3.0,0.8123) (3.2,0.7833) (3.4,0.7738) (3.6,0.7693) (3.8,0.7548) (4.0,0.7108) (4.2,0.6909) (4.4,0.6467) (4.6,0.6383) (4.8,0.6324) (5.0,0.5714) (5.2,0.5843) (5.4,0.5495) (5.6,0.5429) (5.8,0.5162) (6.0,0.4765) (6.2,0.4864) (6.4,0.4582) (6.6,0.4172) (6.8,0.3984) (7.0,0.3773) (7.2,0.3625) (7.4,0.3661) (7.6,0.3284) (7.8,0.2989) (8.0,0)} \\[0.4em]
    \textbf{R} & \histpanel{(-8.0,0.0000) (-7.8,0.0000) (-7.6,0.0000) (-7.4,0.0000) (-7.2,0.0000) (-7.0,0.0000) (-6.8,0.0000) (-6.6,0.0000) (-6.4,0.0001) (-6.2,0.0000) (-6.0,0.0002) (-5.8,0.0001) (-5.6,0.0010) (-5.4,0.0010) (-5.2,0.0017) (-5.0,0.0025) (-4.8,0.0043) (-4.6,0.0060) (-4.4,0.0089) (-4.2,0.0151) (-4.0,0.0237) (-3.8,0.0312) (-3.6,0.0445) (-3.4,0.0629) (-3.2,0.0892) (-3.0,0.1196) (-2.8,0.1591) (-2.6,0.2022) (-2.4,0.2625) (-2.2,0.3245) (-2.0,0.3977) (-1.8,0.4713) (-1.6,0.5594) (-1.4,0.6373) (-1.2,0.7382) (-1.0,0.8258) (-0.8,0.8843) (-0.6,0.9382) (-0.4,0.9670) (-0.2,1.0000) (0.0,0.9905) (0.2,0.9637) (0.4,0.9250) (0.6,0.8881) (0.8,0.8026) (1.0,0.7248) (1.2,0.6562) (1.4,0.5501) (1.6,0.4810) (1.8,0.4096) (2.0,0.3252) (2.2,0.2673) (2.4,0.2070) (2.6,0.1620) (2.8,0.1211) (3.0,0.0890) (3.2,0.0620) (3.4,0.0454) (3.6,0.0343) (3.8,0.0230) (4.0,0.0140) (4.2,0.0108) (4.4,0.0053) (4.6,0.0049) (4.8,0.0017) (5.0,0.0020) (5.2,0.0011) (5.4,0.0010) (5.6,0.0003) (5.8,0.0001) (6.0,0.0001) (6.2,0.0001) (6.4,0.0000) (6.6,0.0000) (6.8,0.0000) (7.0,0.0000) (7.2,0.0000) (7.4,0.0000) (7.6,0.0000) (7.8,0.0000) (8.0,0)} & \histpanel{(-8.0,0.0000) (-7.8,0.0000) (-7.6,0.0000) (-7.4,0.0000) (-7.2,0.0000) (-7.0,0.0001) (-6.8,0.0000) (-6.6,0.0000) (-6.4,0.0003) (-6.2,0.0002) (-6.0,0.0004) (-5.8,0.0004) (-5.6,0.0010) (-5.4,0.0013) (-5.2,0.0024) (-5.0,0.0047) (-4.8,0.0075) (-4.6,0.0102) (-4.4,0.0135) (-4.2,0.0228) (-4.0,0.0314) (-3.8,0.0413) (-3.6,0.0600) (-3.4,0.0777) (-3.2,0.1059) (-3.0,0.1442) (-2.8,0.1880) (-2.6,0.2438) (-2.4,0.2939) (-2.2,0.3605) (-2.0,0.4287) (-1.8,0.5107) (-1.6,0.5769) (-1.4,0.6748) (-1.2,0.7465) (-1.0,0.8360) (-0.8,0.8968) (-0.6,0.9282) (-0.4,0.9790) (-0.2,1.0000) (0.0,0.9960) (0.2,0.9712) (0.4,0.9303) (0.6,0.8618) (0.8,0.8161) (1.0,0.7234) (1.2,0.6382) (1.4,0.5546) (1.6,0.4737) (1.8,0.4029) (2.0,0.3260) (2.2,0.2728) (2.4,0.2051) (2.6,0.1600) (2.8,0.1209) (3.0,0.0894) (3.2,0.0679) (3.4,0.0475) (3.6,0.0336) (3.8,0.0219) (4.0,0.0172) (4.2,0.0118) (4.4,0.0086) (4.6,0.0052) (4.8,0.0033) (5.0,0.0027) (5.2,0.0012) (5.4,0.0011) (5.6,0.0006) (5.8,0.0005) (6.0,0.0001) (6.2,0.0001) (6.4,0.0000) (6.6,0.0000) (6.8,0.0001) (7.0,0.0001) (7.2,0.0000) (7.4,0.0000) (7.6,0.0000) (7.8,0.0000) (8.0,0)} & \histpanel{(-8.0,0.0000) (-7.8,0.0000) (-7.6,0.0000) (-7.4,0.0001) (-7.2,0.0000) (-7.0,0.0000) (-6.8,0.0001) (-6.6,0.0000) (-6.4,0.0003) (-6.2,0.0002) (-6.0,0.0008) (-5.8,0.0010) (-5.6,0.0013) (-5.4,0.0030) (-5.2,0.0037) (-5.0,0.0050) (-4.8,0.0074) (-4.6,0.0110) (-4.4,0.0164) (-4.2,0.0202) (-4.0,0.0320) (-3.8,0.0410) (-3.6,0.0574) (-3.4,0.0756) (-3.2,0.1030) (-3.0,0.1339) (-2.8,0.1800) (-2.6,0.2261) (-2.4,0.2870) (-2.2,0.3525) (-2.0,0.4241) (-1.8,0.4993) (-1.6,0.5979) (-1.4,0.6756) (-1.2,0.7714) (-1.0,0.8382) (-0.8,0.9103) (-0.6,0.9438) (-0.4,0.9863) (-0.2,1.0000) (0.0,0.9916) (0.2,0.9533) (0.4,0.9271) (0.6,0.8647) (0.8,0.8161) (1.0,0.7314) (1.2,0.6583) (1.4,0.5683) (1.6,0.4868) (1.8,0.4253) (2.0,0.3398) (2.2,0.2852) (2.4,0.2178) (2.6,0.1741) (2.8,0.1366) (3.0,0.1041) (3.2,0.0740) (3.4,0.0531) (3.6,0.0397) (3.8,0.0270) (4.0,0.0176) (4.2,0.0118) (4.4,0.0082) (4.6,0.0052) (4.8,0.0030) (5.0,0.0019) (5.2,0.0011) (5.4,0.0007) (5.6,0.0004) (5.8,0.0001) (6.0,0.0001) (6.2,0.0001) (6.4,0.0000) (6.6,0.0000) (6.8,0.0000) (7.0,0.0000) (7.2,0.0000) (7.4,0.0000) (7.6,0.0000) (7.8,0.0000) (8.0,0)} & \histpanel{(-8.0,0.0000) (-7.8,0.0000) (-7.6,0.0000) (-7.4,0.0001) (-7.2,0.0001) (-7.0,0.0000) (-6.8,0.0001) (-6.6,0.0001) (-6.4,0.0003) (-6.2,0.0003) (-6.0,0.0011) (-5.8,0.0012) (-5.6,0.0010) (-5.4,0.0019) (-5.2,0.0030) (-5.0,0.0042) (-4.8,0.0062) (-4.6,0.0101) (-4.4,0.0136) (-4.2,0.0218) (-4.0,0.0266) (-3.8,0.0379) (-3.6,0.0523) (-3.4,0.0698) (-3.2,0.0930) (-3.0,0.1230) (-2.8,0.1558) (-2.6,0.1977) (-2.4,0.2637) (-2.2,0.3338) (-2.0,0.4050) (-1.8,0.4861) (-1.6,0.5879) (-1.4,0.6653) (-1.2,0.7569) (-1.0,0.8493) (-0.8,0.9285) (-0.6,0.9612) (-0.4,0.9932) (-0.2,1.0000) (0.0,0.9963) (0.2,0.9660) (0.4,0.9234) (0.6,0.8554) (0.8,0.7935) (1.0,0.7182) (1.2,0.6268) (1.4,0.5377) (1.6,0.4629) (1.8,0.3834) (2.0,0.3139) (2.2,0.2556) (2.4,0.1982) (2.6,0.1563) (2.8,0.1208) (3.0,0.0938) (3.2,0.0684) (3.4,0.0498) (3.6,0.0375) (3.8,0.0263) (4.0,0.0182) (4.2,0.0126) (4.4,0.0095) (4.6,0.0066) (4.8,0.0045) (5.0,0.0024) (5.2,0.0011) (5.4,0.0014) (5.6,0.0004) (5.8,0.0005) (6.0,0.0001) (6.2,0.0001) (6.4,0.0001) (6.6,0.0001) (6.8,0.0000) (7.0,0.0000) (7.2,0.0000) (7.4,0.0000) (7.6,0.0000) (7.8,0.0000) (8.0,0)} \\[0.4em]
    \textbf{S} & \histpanel{(-8.0,0.0000) (-7.8,0.0000) (-7.6,0.0000) (-7.4,0.0000) (-7.2,0.0000) (-7.0,0.0000) (-6.8,0.0000) (-6.6,0.0000) (-6.4,0.0000) (-6.2,0.0000) (-6.0,0.0000) (-5.8,0.0000) (-5.6,0.0000) (-5.4,0.0000) (-5.2,0.0000) (-5.0,0.0000) (-4.8,0.0000) (-4.6,0.0000) (-4.4,0.0000) (-4.2,0.0000) (-4.0,0.0000) (-3.8,0.0000) (-3.6,0.0000) (-3.4,0.0000) (-3.2,0.0000) (-3.0,0.0000) (-2.8,0.0000) (-2.6,0.0000) (-2.4,0.0000) (-2.2,0.0001) (-2.0,0.0009) (-1.8,0.0036) (-1.6,0.0119) (-1.4,0.0353) (-1.2,0.0910) (-1.0,0.2062) (-0.8,0.3788) (-0.6,0.6189) (-0.4,0.8588) (-0.2,1.0000) (0.0,0.9958) (0.2,0.8527) (0.4,0.6224) (0.6,0.3854) (0.8,0.2042) (1.0,0.0942) (1.2,0.0356) (1.4,0.0118) (1.6,0.0029) (1.8,0.0011) (2.0,0.0001) (2.2,0.0000) (2.4,0.0000) (2.6,0.0000) (2.8,0.0000) (3.0,0.0000) (3.2,0.0000) (3.4,0.0000) (3.6,0.0000) (3.8,0.0000) (4.0,0.0000) (4.2,0.0000) (4.4,0.0000) (4.6,0.0000) (4.8,0.0000) (5.0,0.0000) (5.2,0.0000) (5.4,0.0000) (5.6,0.0000) (5.8,0.0000) (6.0,0.0000) (6.2,0.0000) (6.4,0.0000) (6.6,0.0000) (6.8,0.0000) (7.0,0.0000) (7.2,0.0000) (7.4,0.0000) (7.6,0.0000) (7.8,0.0000) (8.0,0)} & \histpanel{(-8.0,0.0000) (-7.8,0.0000) (-7.6,0.0000) (-7.4,0.0000) (-7.2,0.0000) (-7.0,0.0000) (-6.8,0.0000) (-6.6,0.0000) (-6.4,0.0000) (-6.2,0.0000) (-6.0,0.0000) (-5.8,0.0000) (-5.6,0.0000) (-5.4,0.0000) (-5.2,0.0000) (-5.0,0.0000) (-4.8,0.0000) (-4.6,0.0000) (-4.4,0.0000) (-4.2,0.0000) (-4.0,0.0000) (-3.8,0.0000) (-3.6,0.0000) (-3.4,0.0000) (-3.2,0.0000) (-3.0,0.0000) (-2.8,0.0000) (-2.6,0.0000) (-2.4,0.0000) (-2.2,0.0000) (-2.0,0.0000) (-1.8,0.0000) (-1.6,0.0000) (-1.4,0.0000) (-1.2,0.0000) (-1.0,0.0000) (-0.8,0.0007) (-0.6,0.0229) (-0.4,0.2741) (-0.2,0.9679) (0.0,1.0000) (0.2,0.3095) (0.4,0.0287) (0.6,0.0010) (0.8,0.0000) (1.0,0.0000) (1.2,0.0000) (1.4,0.0000) (1.6,0.0000) (1.8,0.0000) (2.0,0.0000) (2.2,0.0000) (2.4,0.0000) (2.6,0.0000) (2.8,0.0000) (3.0,0.0000) (3.2,0.0000) (3.4,0.0000) (3.6,0.0000) (3.8,0.0000) (4.0,0.0000) (4.2,0.0000) (4.4,0.0000) (4.6,0.0000) (4.8,0.0000) (5.0,0.0000) (5.2,0.0000) (5.4,0.0000) (5.6,0.0000) (5.8,0.0000) (6.0,0.0000) (6.2,0.0000) (6.4,0.0000) (6.6,0.0000) (6.8,0.0000) (7.0,0.0000) (7.2,0.0000) (7.4,0.0000) (7.6,0.0000) (7.8,0.0000) (8.0,0)} & \histpanel{(-8.0,0.0000) (-7.8,0.0000) (-7.6,0.0000) (-7.4,0.0000) (-7.2,0.0000) (-7.0,0.0000) (-6.8,0.0000) (-6.6,0.0000) (-6.4,0.0000) (-6.2,0.0000) (-6.0,0.0000) (-5.8,0.0000) (-5.6,0.0000) (-5.4,0.0000) (-5.2,0.0000) (-5.0,0.0000) (-4.8,0.0000) (-4.6,0.0000) (-4.4,0.0000) (-4.2,0.0000) (-4.0,0.0000) (-3.8,0.0000) (-3.6,0.0000) (-3.4,0.0000) (-3.2,0.0000) (-3.0,0.0000) (-2.8,0.0000) (-2.6,0.0000) (-2.4,0.0000) (-2.2,0.0000) (-2.0,0.0000) (-1.8,0.0000) (-1.6,0.0000) (-1.4,0.0000) (-1.2,0.0000) (-1.0,0.0000) (-0.8,0.0000) (-0.6,0.0000) (-0.4,0.0016) (-0.2,1.0000) (0.0,0.9333) (0.2,0.0023) (0.4,0.0000) (0.6,0.0000) (0.8,0.0000) (1.0,0.0000) (1.2,0.0000) (1.4,0.0000) (1.6,0.0000) (1.8,0.0000) (2.0,0.0000) (2.2,0.0000) (2.4,0.0000) (2.6,0.0000) (2.8,0.0000) (3.0,0.0000) (3.2,0.0000) (3.4,0.0000) (3.6,0.0000) (3.8,0.0000) (4.0,0.0000) (4.2,0.0000) (4.4,0.0000) (4.6,0.0000) (4.8,0.0000) (5.0,0.0000) (5.2,0.0000) (5.4,0.0000) (5.6,0.0000) (5.8,0.0000) (6.0,0.0000) (6.2,0.0000) (6.4,0.0000) (6.6,0.0000) (6.8,0.0000) (7.0,0.0000) (7.2,0.0000) (7.4,0.0000) (7.6,0.0000) (7.8,0.0000) (8.0,0)} & \histpanel{(-8.0,0.0000) (-7.8,0.0000) (-7.6,0.0000) (-7.4,0.0000) (-7.2,0.0000) (-7.0,0.0000) (-6.8,0.0000) (-6.6,0.0000) (-6.4,0.0000) (-6.2,0.0000) (-6.0,0.0000) (-5.8,0.0000) (-5.6,0.0000) (-5.4,0.0000) (-5.2,0.0000) (-5.0,0.0000) (-4.8,0.0000) (-4.6,0.0000) (-4.4,0.0000) (-4.2,0.0000) (-4.0,0.0000) (-3.8,0.0000) (-3.6,0.0000) (-3.4,0.0000) (-3.2,0.0000) (-3.0,0.0000) (-2.8,0.0000) (-2.6,0.0000) (-2.4,0.0000) (-2.2,0.0000) (-2.0,0.0000) (-1.8,0.0000) (-1.6,0.0000) (-1.4,0.0000) (-1.2,0.0000) (-1.0,0.0000) (-0.8,0.0000) (-0.6,0.0000) (-0.4,0.0000) (-0.2,0.9388) (0.0,1.0000) (0.2,0.0000) (0.4,0.0000) (0.6,0.0000) (0.8,0.0000) (1.0,0.0000) (1.2,0.0000) (1.4,0.0000) (1.6,0.0000) (1.8,0.0000) (2.0,0.0000) (2.2,0.0000) (2.4,0.0000) (2.6,0.0000) (2.8,0.0000) (3.0,0.0000) (3.2,0.0000) (3.4,0.0000) (3.6,0.0000) (3.8,0.0000) (4.0,0.0000) (4.2,0.0000) (4.4,0.0000) (4.6,0.0000) (4.8,0.0000) (5.0,0.0000) (5.2,0.0000) (5.4,0.0000) (5.6,0.0000) (5.8,0.0000) (6.0,0.0000) (6.2,0.0000) (6.4,0.0000) (6.6,0.0000) (6.8,0.0000) (7.0,0.0000) (7.2,0.0000) (7.4,0.0000) (7.6,0.0000) (7.8,0.0000) (8.0,0)} \\
    \end{tabular}
    \end{center}

    Associe cada inicialização à linha de histogramas correspondente:
    \textbf{1.} $\sigma_\theta^2 = 1/(4n)$; \quad
    \textbf{2.} $\sigma_\theta^2 = 1/n$; \quad
    \textbf{3.} $\sigma_\theta^2 = 2/n$; \quad
    \textbf{4.} $\sigma_\theta^2 = 4/n$.

    \begin{enumerate}[(a)]
        \item 1-S; 2-R; 3-P; 4-Q
        \item 1-R; 2-Q; 3-S; 4-P
        \item 1-Q; 2-R; 3-P; 4-S
        \item 1-P; 2-Q; 3-S; 4-R
        \item 1-S; 2-P; 3-R; 4-Q
        \item 1-R; 2-S; 3-Q; 4-P
        \item 1-P; 2-S; 3-Q; 4-R
        \item 1-Q; 2-P; 3-R; 4-S
    \end{enumerate}
    %e

\end{enumerate}

\newpage

\subsection*{Gabarito}
\color{blue}
\begin{enumerate}[1)]

    \item c. Com todos os pesos iguais, todas as unidades de uma mesma camada computam a mesma pré-ativação e a mesma saída; no \textit{backward}, recebem gradientes idênticos e são atualizadas de forma idêntica. A simetria nunca é quebrada e a camada se comporta como se tivesse uma única unidade. Note que os gradientes não são nulos (a): o problema não é ausência de atualização, e sim atualização idêntica em todas as unidades.

    \item b. Pesos com variância grande produzem pré-ativações de magnitude grande, que colocam a sigmoid nas regiões saturadas onde a derivada é praticamente zero; no \textit{backpropagation} isso dissipa o gradiente antes de chegar às camadas iniciais. Unidades mortas (a) são um fenômeno da ReLU, não da sigmoid, e (c) descreve o cenário oposto, de $\sigma$ pequeno demais.

    \item \textit{Vanishing Gradient}: a sigmoid satura para entradas de magnitude grande e sua derivada vale no máximo $0.25$; ao multiplicar essas derivadas através das camadas no \textit{backpropagation}, o gradiente encolhe exponencialmente e as primeiras camadas praticamente param de aprender.

    \item Unidades mortas (\textit{dead units}): se os pesos e o bias de uma unidade a levam a produzir pré-ativação negativa para todas as entradas do conjunto, a saída da ReLU é sempre zero e sua derivada também; o gradiente que passa pela unidade é nulo e ela nunca mais é atualizada.

    \item Inicializar os pesos da rede neural de modo que não ocorra explosão nem dissipação de gradiente. A ideia central é escolher a variância dos pesos de forma que a variância das pré-ativações (e, por consequência, dos gradientes) se preserve de camada para camada: a He usa $\sigma_\theta^2 = 2/n^{(l-1)}$ para compensar o fator $\frac{1}{2}$ que a ReLU introduz no segundo momento, e a Xavier usa $\sigma_\theta^2 = 1/n^{(l-1)}$ para ativações aproximadamente lineares perto da origem, como a tanh.

    \item Pela definição da ReLU,
    \begin{equation*}
        E[\mathrm{ReLU}(a)^2]=\int_{-\infty}^{\infty}\max(0,a)^2\,p(a)\,da=\int_{0}^{\infty}a^2\,p(a)\,da.
    \end{equation*}
    Pela simetria de $p$ em torno de zero, $\int_{-\infty}^{0}a^2\,p(a)\,da=\int_{0}^{\infty}a^2\,p(a)\,da$, então $E[a^2]=2\int_{0}^{\infty}a^2\,p(a)\,da$. Como $E[a]=0$, temos $E[a^2]=\mathrm{Var}[a]=\sigma^2$, e portanto $E[\mathrm{ReLU}(a)^2]=\frac{E[a^2]}{2}=\frac{\sigma^2}{2}$.

    \item
    \begin{enumerate}[a)]
        \item Pela linearidade da esperança, $E[z^{(l)}_i]=E[b^{(l)}_i]+\sum_j E[\theta^{(l)}_{ij}\,\varphi(z^{(l-1)}_j)]$. O bias é inicializado em zero e, pela independência entre pesos e ativações, cada termo do somatório fatora em $E[\theta^{(l)}_{ij}]\,E[\varphi(z^{(l-1)}_j)]=0\cdot E[\varphi(z^{(l-1)}_j)]=0$. Logo $E[z^{(l)}_i]=0$.
        \item Como $E[z^{(l)}_i]=0$, vale $\sigma^2_{z^{(l)}}=E[(z^{(l)}_i)^2]$. Ao expandir o quadrado do somatório, os termos cruzados desaparecem pela independência dos pesos e por $E[\theta]=0$, restando $\sigma^2_{z^{(l)}}=\sum_j E[(\theta^{(l)}_{ij})^2]\,E[\varphi(z^{(l-1)}_j)^2]=\sigma_\theta^2\sum_j E[\varphi(z^{(l-1)}_j)^2]$. Pelo exercício anterior, $E[\mathrm{ReLU}(z^{(l-1)}_j)^2]=\sigma^2_{z^{(l-1)}}/2$ para cada uma das $n^{(l-1)}$ unidades, o que dá $\sigma^2_{z^{(l)}} = \frac{n^{(l-1)}}{2}\,\sigma_\theta^2\,\sigma^2_{z^{(l-1)}}$.
        \item Impondo $\sigma^2_{z^{(l)}}=\sigma^2_{z^{(l-1)}}$ na recorrência, a constante multiplicativa deve valer 1: $\frac{n^{(l-1)}}{2}\,\sigma_\theta^2=1$, ou seja, $\sigma_\theta^2=\frac{2}{n^{(l-1)}}$, que é a inicialização \textit{He}.
        \item Com $\varphi(z)\approx z$, o segundo momento da ativação passa a ser $E[\varphi(z)^2]\approx\sigma^2_{z^{(l-1)}}$, sem o fator $\frac{1}{2}$ da ReLU. A recorrência vira $\sigma^2_{z^{(l)}}=n^{(l-1)}\,\sigma_\theta^2\,\sigma^2_{z^{(l-1)}}$ e a condição de preservação dá $\sigma_\theta^2=\frac{1}{n^{(l-1)}}$, que é a inicialização \textit{Xavier}.
    \end{enumerate}

    \item d. He propôs (\textit{He et al., 2015}) usar $\mathrm{Var}(W) = 2/n_{\text{in}}$ assumindo ReLU; o fator 2 compensa a probabilidade aproximada de $1/2$ de a pré-ativação ser zerada pela ReLU, mantendo a variância das ativações estável ao longo da profundidade. Em (a), pesos zero geram unidades idênticas e o gradiente não quebra essa simetria; em (b), a mesma constante para todos os pesos também viola a quebra de simetria entre unidades; em (c), Xavier é apropriada para tanh/sigmoide (regime aproximadamente linear), e para ReLU o fator 2 ausente faz com que a variância das ativações encolha camada a camada; em (e), uniforme em $[-1, 1]$ ignora $n_{\text{in}}$, e em redes profundas isso faz a variância das ativações explodir ou colapsar; em (f), PCA produz uma inicialização determinística, que não é prática padrão e não tem garantia teórica para esse fim; em (g), apenas inicializar os \textit{biases} aleatoriamente não basta, pois pesos com mesmo valor produzem saídas idênticas independentemente do \textit{bias}; em (h), $\mathcal{N}(0, 1)$ sem reescala em rede profunda sob ReLU faz a variância explodir, e não basta diminuir a taxa de aprendizado.

    \item e. Com vieses nulos, a análise vista em aula dá, na primeira camada (a entrada não passa por ReLU e, como $\mathbb{E}[x_j] = 0$, o segundo momento é $\mathbb{E}[x_j^2] = \mathrm{Var}[x_j] = 1$),
    \[
    \sigma^2_{z^{(1)}} = n\,\sigma_\theta^2 ,
    \]
    e, a partir da segunda camada, $\sigma^2_{z^{(l)}} = \tfrac{n}{2}\,\sigma_\theta^2\,\sigma^2_{z^{(l-1)}}$: a ReLU zera metade da distribuição e o segundo momento cai pela metade. Assim, o desvio padrão da primeira camada é $\sqrt{n\sigma_\theta^2}$ e, a cada camada seguinte, ele é multiplicado por $\sqrt{n\sigma_\theta^2/2}$:
    \[
    \begin{array}{c|c|c|cccc|c}
     & n\sigma_\theta^2 & \text{fator por camada} & \mathbf{Z}^{(1)} & \mathbf{Z}^{(2)} & \mathbf{Z}^{(3)} & \mathbf{Z}^{(4)} & \text{linha}\\\hline
    1/(4n) & 0{,}25 & \sqrt{0{,}125} \approx 0{,}35 & 0{,}50 & 0{,}18 & 0{,}06 & 0{,}02 & \text{S}\\
    1/n & 1 & \sqrt{0{,}5} \approx 0{,}71 & 1{,}00 & 0{,}71 & 0{,}50 & 0{,}35 & \text{P}\\
    2/n & 2 & 1 & 1{,}41 & 1{,}41 & 1{,}41 & 1{,}41 & \text{R}\\
    4/n & 4 & \sqrt{2} \approx 1{,}41 & 2{,}00 & 2{,}83 & 4{,}00 & 5{,}66 & \text{Q}
    \end{array}
    \]
    Os desvios padrão medidos na simulação que gerou a figura ficam próximos desses valores (o maior desvio é em Q, $\mathbf{Z}^{(4)}$: $5{,}54$ medido contra $5{,}66$ previsto). A leitura da figura: R mantém a mesma largura nas quatro camadas (inicialização de He); Q se alarga a cada camada (\textit{exploding}); P e S encolhem, mas S já começa mais estreita que P em $\mathbf{Z}^{(1)}$ e vira uma agulha em $\mathbf{Z}^{(2)}$, enquanto P ainda tem largura visível em $\mathbf{Z}^{(4)}$. Logo, 1-S, 2-P, 3-R e 4-Q. Em (a), $1/n$ e $2/n$ foram trocados, como se $\sigma_\theta^2 = 1/n$ (Xavier com \textit{fan-in}) preservasse a variância; isso vale para a tanh perto da origem, mas a ReLU exige o fator $2$. (b) troca os dois pares vizinhos, $1/(4n)$ com $1/n$ e $2/n$ com $4/n$, ignorando que S já é a mais estreita na primeira camada e tomando a linha que se alarga pela estável. (c) inverte toda a ordem, como se pesos com variância maior produzissem pré-ativações mais estreitas. (d) coloca $1/n$ na linha que explode e $2/n$ na que colapsa, deixando $4/n$ na estável. (f) coloca $1/(4n)$ na linha estável e $4/n$ em P, que encolhe a cada camada. (g) troca apenas os extremos, $1/(4n)$ e $4/n$. (h) coloca $1/(4n)$ na linha estável, $1/n$ na que explode, $2/n$ na que colapsa e $4/n$ em P, que encolhe.
\end{enumerate}

\end{document}
